% Source: https://artofproblemsolving.com/community/c3h1988962_about_pi
% Stable identity: aops-topic-1988962-q01
(a) Polynomial long division gives
\[
x^4(1-x)^4=x^8-4x^7+6x^6-4x^5+x^4=(1+x^2)\left(x^6-4x^5+5x^4-4x^2+4\right)-4 ,
\]
so
\[
\int_0^1\dfrac{x^4(1-x)^4}{1+x^2}\,dx=\int_0^1\left(x^6-4x^5+5x^4-4x^2+4\right)dx-\int_0^1\dfrac{4}{1+x^2}\,dx .
\]
The first integral is $\dfrac{1}{7}-\dfrac{4}{6}+1-\dfrac{4}{3}+4=\dfrac{22}{7}$, and the second is $4\arctan1=\pi$. The integrand is positive on $(0,1)$, so the integral is positive and $\dfrac{22}{7}-\pi>0$.

(b) On $[0,1]$ we have $1\le1+x^2\le2$, with equality only at the ends, so
\[
\dfrac{1}{2}\int_0^1x^4(1-x)^4\,dx<\dfrac{22}{7}-\pi<\int_0^1x^4(1-x)^4\,dx .
\]
Expanding, $\int_0^1x^4(1-x)^4\,dx=\dfrac{1}{5}-\dfrac{4}{6}+\dfrac{6}{7}-\dfrac{4}{8}+\dfrac{1}{9}=\dfrac{1}{630}$. Hence
\[
\dfrac{1}{1260}<\dfrac{22}{7}-\pi<\dfrac{1}{630}. \qquad\blacksquare
\]
So $\dfrac{22}{7}$ is close to $\pi$ but not equal to it. In fact $\pi$ is irrational, so no fraction of whole numbers equals it; $\dfrac{22}{7}$ is simply a convenient approximation, and other fractions such as $\dfrac{355}{113}$ are better ones.
